
\chapter{语法知识}


\section*{一、词类总表}

\begin{center}
\renewcommand{\arraystretch}{1.6}
\setlength{\tabcolsep}{6pt}

\subsection*{实词（表示实在意义，能独立成句或作句子成分）}

\begin{longtable}{|>{\centering\arraybackslash}p{2.2cm}|>{\centering\arraybackslash}p{2.8cm}|>{\arraybackslash}p{5.8cm}|>{\arraybackslash}p{4.5cm}|}
\hline
\rowcolor{titlebg}\textbf{词类} & \textbf{定义} & \textbf{分类举例} & \highlight{语法特点} \\
\hline
\endfirsthead
\hline
\rowcolor{titlebg}\textbf{词类} & \textbf{定义} & \textbf{分类举例} & \highlight{语法特点} \\
\hline
\endhead

名词 & 表示人和事物名称 & 人：同志、作家；事物：河流、高山；抽象：政治、科学；时间：上午、夏天；处所：上海、中国；方位：上、下 & \highlight{①可加“们”表多数；②方位词可组方位短语；③一般不受副词修饰} \\
\hline
动词 & 表示动作行为、发展变化、心理活动等 & 动作：坐、听；存现：有、发生；心理：爱、恨；使令：叫、让；能愿：能、会；趋向：来、去；判断：是 & \highlight{①可受“不”修饰；②可带“着、了、过”；③部分可重叠；④“是”联结主宾；⑤能愿动词后不能跟名词；⑥趋向动词可作补语} \\
\hline
形容词 & 表示事物的形状、性质、状态 & 形状：高、矮；性质：漂亮、结实；状态：快、慢 & \highlight{①部分可重叠加强语义；②大多可受“很”修饰} \\
\hline
数词 & 表示数目 & 确数（分数、整数、倍数）；概数：几、许多；序数：第一、老三 & \highlight{①增加可用分数或倍数；②减少只能用分数} \\
\hline
量词 & 表示事物和动作、行为的单位 & 物量词：个、只；动量词：次、回；借用名词：脚、年 & \highlight{①常与数词连用组数量短语；②物量短语用在名词前；③动量短语用在动词后} \\
\hline
代词 & 起代替或指示作用 & 人称代词、疑问代词、指示代词 & \highlight{①“您”不用于复数；②“他们”可兼指，“她们”专指女性；③“我们”与“咱们”有别；④“那”远指“这”近指；⑤指代不明造成病句} \\
\hline
\end{longtable}

\vspace{0.5cm}

\subsection*{虚词（一般不表示实在意义，主要表示语法关系）}

\begin{longtable}{|>{\centering\arraybackslash}p{2.2cm}|>{\centering\arraybackslash}p{2.8cm}|>{\arraybackslash}p{5.8cm}|>{\arraybackslash}p{4.5cm}|}
\hline
\rowcolor{titlebg}\textbf{词类} & \textbf{定义} & \textbf{分类举例} & \highlight{语法特点} \\
\hline
\endfirsthead
\hline
\rowcolor{titlebg}\textbf{词类} & \textbf{定义} & \textbf{分类举例} & \highlight{语法特点} \\
\hline
\endhead

副词 & 用在动词、形容词前，表示程度、范围、时间、频率、情势、语气等 & 范围：都、全；语气：可、倒；否定：不、没；时间：刚、恰好；程度：很、极；情势：仿佛、渐渐 & \highlight{①修饰限制动词或形容词，作状语；②有时作补语；③不能修饰名词、代词} \\
\hline
连词 & 连接词、短语或句子 & 一般连词：和、与、并、或、及；关联词：不但…而且、虽然…但是 & \highlight{①一般连词前后可调换；②关联词用于复句} \\
\hline
介词 & 用在名词、代词前，表示起止、方向、处所、时间、对象、方式、原因、目的、比较等 & 自、从、以、当、为、按照、由于、对于、为了、到、和、跟、把、比、在、关于、除了、同、对、向、往、朝 & \highlight{①组成介宾短语；②作状语或补语；③作定语时后必须带“的”} \\
\hline
助词 & 附着在实词、短语或句子后面起辅助作用 & 结构助词：的、地、得；动态助词：着、了、过；语气助词：吗、吧、呢 & 附着于其他成分，辅助表达 \\
\hline
叹词 & 表示感叹、呼唤、应答等声音 & 啊、嗯等 & \highlight{一般独立成句，用逗号或感叹号隔开} \\
\hline
拟声词 & 摹拟人或事物的声音 & 如：哗啦、叮当、轰隆 & 在句子中相当于一个形容词 \\
\hline
\end{longtable}
\end{center}

\section*{二、词类辨别要诀（重点标注）}

\begin{center}
\renewcommand{\arraystretch}{1.8}
\begin{tabular}{|>{\centering\arraybackslash}p{2.8cm}|>{\arraybackslash}p{11cm}|}
\hline
\rowcolor{titlebg}\textbf{易混词类} & \highlight{辨别方法} \\
\hline
名词 vs 非名词 & 名词前 \highlight{不能} 加“不”和“很” \\
\hline
形容词 vs 动词 & 形容词 \highlight{可以} 用“很”来修饰，动词前 \highlight{不能} 加“很”（心理动词除外） \\
\hline
形容词 vs 副词 & 形容词能修饰名词，前面能加“很”；副词 \highlight{不能} 修饰名词，前面 \highlight{不能} 加“很” \\
\hline
连词 vs 介词 & 前后 \highlight{能互换} 的是连词，前后 \highlight{不能互换} 的是介词 \\
\hline
动词 vs 介词 & 作谓语中心语的是动词，组成介宾短语修饰、补充动形的是介词 \\
\hline
语气助词 vs 叹词 & 语气助词一般在句尾，叹词 \highlight{独立成句}，一般在句首 \\
\hline
介词 vs 副词 & 介词后面跟名词、代词；副词后面是动词或形容词 \\
\hline
\end{tabular}
\end{center}

\section*{三、短语类型总览}

\begin{center}
\renewcommand{\arraystretch}{1.6}
\setlength{\tabcolsep}{5pt}
\begin{longtable}{|>{\centering\arraybackslash}p{2.5cm}|>{\arraybackslash}p{3.5cm}|>{\arraybackslash}p{8.5cm}|}
\hline
\rowcolor{titlebg}\textbf{短语类型} & \textbf{基本结构} & \highlight{特点或语法功能} \\
\hline
\endfirsthead
\hline
\rowcolor{titlebg}\textbf{短语类型} & \textbf{基本结构} & \highlight{特点或语法功能} \\
\hline
\endhead

并列短语 & 名+名、动+动、形+形、数量+数量 & 词性一致，平等关系，一般可颠倒，可直接组合也可用虚词 \\
\hline
偏正短语 & \begin{tabular}{@{}l@{}}定中：形+名、数量+名\\ 状中：副+动、副+形\end{tabular} & 修饰限制关系，定语用（ ）表示，状语用〔 〕表示 \\
\hline
动宾短语 & 动+名、动+代 & 支配与被支配关系，宾语在动词后回答“谁”“什么” \\
\hline
补充短语 & 动/形+补语 & 补语用＜ ＞表示，补充说明怎么样、多久、多少，常用“得”连接 \\
\hline
主谓短语 & 名(代)+动/形；名+名（今天星期一） & 陈述与被陈述关系，加上语气和标点即为完整单句 \\
\hline
介宾短语 & 介词+名/代 & 整体作句子成分，常作状语或补语，作定语时要带“的” \\
\hline
“的”字短语 & 动/形+的、动宾短语+的 & \highlight{相当于名词}，常作主语或宾语 \\
\hline
\end{longtable}
\end{center}

\section*{四、副词分类速查}

\begin{center}
\renewcommand{\arraystretch}{1.5}
\begin{tabular}{|>{\centering\arraybackslash}p{2.5cm}|>{\centering\arraybackslash}p{4.5cm}|>{\centering\arraybackslash}p{4.5cm}|}
\hline
\rowcolor{titlebg}\textbf{类别} & \textbf{常见副词} & \textbf{示例} \\
\hline
范围 & 都、全、只、一共、统统 & 我们\highlight{都}去了 \\
\hline
语气 & 可、倒、难道、究竟、偏偏 & 他\highlight{倒}是来了 \\
\hline
否定 & 不、没、没有、别 & \highlight{不}知道 \\
\hline
时间 & 刚、恰好、已经、正在、马上 & 他\highlight{刚}走 \\
\hline
程度 & 很、极、非常、十分、太、更 & \highlight{很}好 \\
\hline
情势 & 仿佛、渐渐、忽然、猛然 & 天\highlight{渐渐}暗了 \\
\hline
\end{tabular}
\end{center}